% Chapter 5

\chapter{进一步研究}

\section{预期共识定价效应的机制分析}

本节将通过实验金融的方式，分析预期共识定价效应的形成机制。
根据Miller的假说\cite{miller1977risk}，当投资者面临卖空约束时，
悲观投资者无法通过卖空交易充分表达其负面预期，资产价格将主要由乐观投资者的信念决定；
投资者预期分歧越大，预期共识越低，股价被乐观预期高估的程度就越严重。
唐国豪等学者\cite{唐国豪2025投资者预期协同}
使用多个神经网络模型，模拟异质性个体投资者的预期，
并且通过融券可行性和机构投资者持股比例，
来检验融券可行性在预期共识定价效应中的作用。
研究发现，融券可行的证券，其预期共识定价效应不显著，
而融券不可行的证券，其预期共识定价效应显著。
该研究表明，融券可行性是预期共识定价效应的重要形成机制。

本文将使用实验金融的方式来研究融券可行性在预期共识定价效应中的作用。
具体来说，在基准模型中，约定了不允许融资融券，即对于任何一个神经网络的输出，
进行截断，使得输出为非负数。

\subsection{放松融资融券约束}
在本节，对该约束进行放松，允许Agent进行融资和融券。
具体来说，对于最终决策：

\[
\vec{action}_{a,\vec{portfolio},t} = (w_1, w_2, \ldots, w_{n+1})
\]

决策需要满足约束：

\textbf{1)} 权重之和为1，即$\sum_{i=1}^{n+1} w_i = 1$；  

\textbf{2)} 允许$w_i$为负数，当$w_i<0, i \neq n+1$时，
表示做空证券$i$，获取总资金的$|w_i|$倍的资金；
如果$w_{n+1} < 0$，表示借入资金购买证券，获取总资金的$|w_{n+1}|$倍的资金；
为了避免无限制的融资融券，设定一个最大杠杆率$L_{max} \geq 0$，使得：
\[
\sum_{i=1}^{n+1} |w_i| \cdot I(w_{i} < 0) \leq L_{max}
\]

其中，$I(x)$为指示函数，当$x$为真时，取1，否则取0。

由于使用Tanh作为激活函数，输出$o_{i}$范围为$(-1, 1)$，
因此，为了满足上述约束，需要对$o_{i}$进行归一化处理，具体来说：

\begin{enumerate}
    \item 计算所有小于0的$o_{i}$的绝对值之和$S$，即：
    \[
    S = \sum_{i=1}^{n} |o_i| \cdot I(o_i < 0)
    \]
    如果$S > L_{max}$，则按照$\frac{L_{max}}{S}$的比例，缩小所有小于0的$o_{i}$：
    \[
    o_{i}^{norm} = o_{i} \cdot \frac{L_{max}}{S} \quad if \quad o_{i} < 0
    \]
    否则，不变，即：
    \[
    o_{i}^{norm} = o_{i} \quad if \quad o_{i} < 0
    \]

    \item 将所有大于0的$o_{i}$归一化，使得其和为$1 + L_{max}$，即：
    \[
    o_{i}^{norm} = o_{i} \cdot \frac{1 + L_{max}}{\sum_{i=1}^{n} o_{i}} \quad if \quad o_{i} \geq 0
    \]

    \item 归一化后的输出记作：
    \[
    \textbf{o}_{a,\vec{portfolio},t}^{norm} = (o_{1}^{norm}, o_{2}^{norm}, \ldots, o_{n+1}^{norm})
    \]

    \item 归一化后的输出就是最终的决策向量，记作：
    \[
    \vec{action}_{a,\vec{portfolio},t} = (w_1, w_2, \ldots, w_{n+1})
    \]
    其中，$w_i = o_{i}^{norm}$。
\end{enumerate}

\subsection{实验结果}
分别设定融资融券约束$L_{max} = 0.25, 0.5$，
再进行实验，对于结果进行分组检验和Fama-MacBeth回归。
由于篇幅限制，仅展示对冲组合的平均收益、夏普比率、FF5-$\alpha$与5控制变量FM回归结果。
结果如表~\ref{tab:mechanism_analysis_margin}所示。

{\small
\begin{table}[H]
  \centering
  \caption{机制分析}
  \label{tab:mechanism_analysis_margin}
  \makebox[0pt][c]{
  \begin{tabular*}{0.95\textwidth}{@{\extracolsep{\fill}}ccccc}
      \toprule
      融资融券上限 & 平均收益 & 夏普比率 & FF5 Alpha & FM回归 \\
      \midrule
      0 & \makecell{0.380$^{***}$ \\ \text{[4.045]} \\ \text{(0.000)}} & 26.840 & \makecell{0.207$^{***}$ \\ \text{[2.247]} \\ \text{(0.000)}} & \makecell{0.139$^{***}$ \\ \text{[14.440]} \\ \text{(0.000)}} \\
      \midrule
      0.25 & \makecell{0.063$^{*}$ \\ \text{[1.120]} \\ \text{(0.064)}} & 12.562 & \makecell{0.083$^{*}$ \\ \text{[0.949]} \\ \text{(0.072)}} & \makecell{0.013$^{**}$ \\ \text{[1.570]} \\ \text{(0.035)}} \\
      \midrule
      0.5 & \makecell{0.001 \\ \text{[0.312]} \\ \text{(0.297)}} & 6.248 & \makecell{$0.002$ \\ \text{[0.412]} \\ \text{(0.424)}} & \makecell{0.063$^{*}$ \\ \text{[1.002]} \\ \text{(0.051)}} \\
      \bottomrule
    \end{tabular*}
  }
\end{table}
}

 \FloatBarrier

从表中结果可见，融资融券约束松紧与预期共识定价效应直接相关，
约束越宽松，定价效应越弱。融
资融券上限为0的严格约束下，对冲组合平均收益0.380、FF5-α及FM回归系数均在1\%水平高度显著，
夏普比率达26.840\%，定价效应极强；
约束放松至0.25时，组合收益大幅回落，平均收益仅0.063（10\%水平边际显著），
夏普比率降至12.562，各项指标显著性同步弱化；约束进一步放宽至0.5后，
超额收益近乎消失，平均收益无统计显著性，FF5-α不再显著，
仅FM回归系数勉强维持10\%水平显著。整体而言，随着融资融券限制放松，组合收益、夏普比率及回归显著性均大幅下滑，
说明约束收紧时投资者难以充分交易修正预期，催生预期共识定价效应，约束放松则该效应显著弱化，
由此证实融资融券限制是预期共识定价效应的重要形成机制。

随着融资融券限制的放松，
对冲组合的平均收益和夏普比率明显下降，FF5-$\alpha$和FM回归的显著性也明显下降。
该结果表明，随着融资融券的放松，投资者可以更容易的表达其预期，
从而修正了预期共识的定价效应，导致预期共识的定价效应不显著。
因此，融资融券限制是预期共识定价效应的重要形成机制。


\section{异质性分析}
按照市值、企业产权性质和股市周期三个维度，检验结论的异质性。
具体来说：
第一，每一个月份，对企业的市值进行排序，按照是否大于中位数，划分为大市值与小市值两组；
第二，按照是否为国有企业，划分为非国有企业与国有企业两组；
第三，参考Maheu和McCurdy的方法\cite{maheu2000identifying}，使用沪深300指数，划分每一个月份为牛市或熊市周期。
其具体做法是，将股市变动视作一个马尔可夫过程，使用最大似然估计方法，估计出牛市和熊市状态的转移概率矩阵，
并根据转移概率矩阵，划分每一个月份为牛市或熊市周期。

在按照异质性分组后，每一组内部，
再将观测按各自预期共识水平划分为5个分组，
保持$2 \times 5 = 10$个分组。
在每个分组内部，分别进行分组检验和Fama-MacBeth回归。

由于篇幅限制，仅展示对冲组合的平均收益、夏普比率、FF5-$\alpha$与5控制变量FM回归结果。


\subsection{基于企业市值的异质性分析}
为考察预期共识效应在不同市值股票中的表现差异，本文将样本按市值高低划分为大市值与小市值两组，
分别进行分组检验与Fama-MacBeth回归，结果如表~\ref{tab:big_small}所示。

{\small
\begin{table}[H]
  \centering
  \caption{基于企业市值的异质性分析}
  \label{tab:big_small}
  \makebox[0pt][c]{
  \begin{tabular*}{0.95\textwidth}{@{\extracolsep{\fill}}ccccc}
      \toprule
      市值 & 平均收益 & 夏普比率 & FF5 Alpha & FM回归 \\
      \midrule
      大市值 & \makecell{0.589$^{***}$ \\ \text{[4.913]} \\ \text{(0.000)}} & 33.120 & \makecell{0.406$^{***}$ \\ \text{[3.427]} \\ \text{(0.001)}} & \makecell{0.126$^{***}$ \\ \text{[14.450]} \\ \text{(0.001)}}  \\
      \midrule
      小市值 & \makecell{1.216$^{***}$ \\ \text{[9.613]} \\ \text{(0.000)}} & 98.380 & \makecell{1.094$^{***}$ \\ \text{[8.686]} \\ \text{(0.000)}} & \makecell{0.154$^{***}$ \\ \text{[12.270]} \\ \text{(0.000)}} \\
      \bottomrule
    \end{tabular*}
  }
\end{table}
}

 \FloatBarrier

从分组检验结果来看，预期共识效应在大市值与小市值样本中均显著成立。
在大市值样本中，平均收益率随组合呈现单调递增趋势，
对冲组合的平均收益率为0.589\%，在 1\% 水平上显著，
对应的 CAPM-$\alpha$、FF3-$\alpha$ 与 FF5-$\alpha$ 均在 1\% 水平上显著为正。

在小市值样本中，预期共识的定价效应更为突出：
对冲组合的平均收益率高达1.216\%,在 1\% 水平上显著，
对应的 CAPM-$\alpha$、FF3-$\alpha$ 与 FF5-$\alpha$ 均在 1\% 水平上显著为正，
且其数值远超大市值样本。
同时，小市值对冲组合的夏普比率为 98.380\%，
也显著高于大市值样本的33.120\%，表明小市值中预期共识策略的收益风险比更优。

Fama-MacBeth回归结果进一步验证了上述结论。
在大市值样本中，无论是否加入控制变量，
预期共识指标MA的回归系数均在1\% 水平上显著为正：
在加入5个控制变量后，MA 系数仍为0.126。
在小市值样本中，MA的回归系数同样在1\% 水平上显著为正，且数值明显更大：
加入5个控制变量后为0.154，表明预期共识对小市值股票收益的边际影响更强。

上述异质性结果可以从错误定价视角加以解释，
核心逻辑为：小市值股票面临更严格的卖空约束和更严重的有限套利限制，
使其错误定价程度更深，进而放大了预期共识的定价效应。
大市值股票多为两融核心标的，卖空约束弱，价格高估空间有限；
小市值股票普遍两融覆盖不足、券源稀缺，卖空约束极强，
低共识下悲观预期无法入市导致股价显著高估，高共识则对应理性定价，
因此小市值高低共识组合的收益差异更突出。

\subsection{股权性质异质性分析}
为考察预期共识效应在不同股权性质企业中的表现差异，
本文将样本划分为非国有企业与国有企业两组，
分别进行分组检验与Fama-MacBeth回归，
结果如表~\ref{tab:ownership}所示。

{\small
\begin{table}[H]
  \centering
  \caption{股权性质异质性分析}
  \label{tab:ownership}
  \makebox[0pt][c]{
  \begin{tabular*}{0.95\textwidth}{@{\extracolsep{\fill}}ccccc}
      \toprule
      股权性质 & 平均收益 & 夏普比率 & FF5 Alpha & FM回归 \\
      \midrule
      国有企业 & \makecell{0.313 \\ \text{[1.627]} \\ \text{(0.104)}} & 11.110 & \makecell{0.128 \\ \text{[0.714]} \\ \text{(0.476)}} & \makecell{0.111$^{***}$ \\ \text{[14.900]} \\ \text{(0.000)}} \\
      \midrule
      非国有企业 & \makecell{0.907$^{***}$ \\ \text{[7.326]} \\ \text{(0.000)}} & 55.480 & \makecell{0.711$^{***}$ \\ \text{[5.869]} \\ \text{(0.000)}} & \makecell{0.151$^{***}$ \\ \text{[12.700]} \\ \text{(0.000)}} \\
      \bottomrule
    \end{tabular*}
  }
\end{table}
}

从分组检验结果来看，预期共识效应在非国有企业中表现显著，而在国有企业中明显弱化。
在非国有企业样本中，平均收益率随组合呈现单调递增趋势：
对冲组合的平均收益率高达0.907\%，在1\%水平上显著，
对应的$\alpha$均在1\%水平上显著为正。
在国有企业样本中，尽管组合收益率从低到高有一定递增趋势，
但对冲组合的平均收益率仅为0.313\%，且统计上不显著；对应的$\alpha$也均不显著，
表明国有企业中预期共识的定价效应在分组检验层面未得到充分体现。

Fama-MacBeth回归结果进一步验证了上述结论。
在非国有企业样本中，
无论是否加入控制变量，预期共识指标MA的回归系数均在1\%水平上显著为正：
在加入5个控制变量后，MA系数仍为0.151。
在国有企业样本中，MA的回归系数虽仍在1\%水平上显著为正，
但数值明显更小：加入5个控制变量后为0.111，
表明预期共识对国有企业股票收益的边际影响更弱。

上述异质性结果可以从错误定价视角并结合国企与非国企的制度差异加以解释，
核心逻辑为：非国企面临更严重的信息不对称且缺乏政府隐性背书，
使其错误定价程度更深，进而放大了预期共识的定价效应。
一方面，国企受监管更严格，信息披露规范、分析师覆盖充分，
投资者信息不对称程度较低，预期共识的增量定价效应有限；
非国企信息透明度偏低，预期共识作为市场对公司价值的认可信号，其“定价锚” 作用显著更强。
另一方面，国企的政府信用背书带来股价估值刚性，
天然压缩了错误定价空间；非国企缺乏此类背书，估值高度依赖市场对其经营成长的预期，
分歧加剧时极易出现显著错误定价，进一步强化了预期共识的定价效应。

\subsection{市场周期异质性分析}
为考察预期共识效应在不同市场周期中的表现差异，
分别进行分组检验与Fama-MacBeth回归，结果如表~\ref{tab:cycle}所示。

{\small
\begin{table}[H]
  \centering
  \caption{市场周期异质性分析}
  \label{tab:cycle}
  \makebox[0pt][c]{
  \begin{tabular*}{0.95\textwidth}{@{\extracolsep{\fill}}ccccc}
      \toprule
      市场周期 & 平均收益 & 夏普比率 & FF5 Alpha & FM回归 \\
      \midrule
      熊市 & \makecell{0.688 \\ {[1.548]} \\ {(0.122)}} & 19.710 & \makecell{0.195 \\ {[0.461]} \\ {(0.645)}} & \makecell{0.155$^{***}$ \\ {[9.924]} \\ {(0.000)}} \\
      \midrule
      牛市 & \makecell{0.511$^{***}$ \\ {[3.997]} \\ {(0.000)}} & 29.640 & \makecell{0.364$^{**}$ \\ {[2.581]} \\ {(0.010)}} & \makecell{0.132$^{***}$ \\ {[11.190]} \\ {(0.000)}} \\
      \bottomrule
    \end{tabular*}
  }
\end{table}
}

从分组检验结果来看，预期共识效应在牛市中表现显著，而在熊市中分组检验的显著性有所下降。
在牛市样本中，平均收益率随组合呈现单调递增趋势：
对冲组合（做多高预期共识组合、做空低预期共识组合）的平均收益率为0.511\%，在1\%水平上显著，
对应的$\alpha$均在1\%或5\%水平上显著为正。
在熊市样本中，尽管组合收益率从低到高有一定递增趋势，
对冲组合的平均收益率为0.688\%，但统计上不显著；
对应的$\alpha$也均不显著。

Fama-MacBeth回归结果进一步验证了预期共识效应在不同市场周期中的存在性。
在牛市样本中，无论是否加入控制变量，预期共识指标MA的回归系数均在1\%水平上显著为正：
在加入5个控制变量后，MA系数仍为0.132。
在熊市样本中，尽管仍然有显著的正向系数，但显著性有所下降。

上述异质性结果可以基于行为金融学中投资者情绪理论与错误定价理论加以解释，
核心逻辑为：牛市中高涨的投资者情绪与卖空约束放大了错误定价，熊市中预期共识则发挥了更强的“安全信号”作用。
具体而言：第一，牛市中投资者乐观情绪主导市场，低预期共识对应乐观投资者的极端预期，
在卖空约束下悲观预期无法充分反映于价格，导致股价显著高估；
高预期共识则对应更理性的定价，因此高低共识组合的收益差异显著。
第二，熊市中投资者风险厌恶程度大幅提升，预期共识成为验证公司价值与抗风险能力的强安全信号，
而预期分歧大的股票更易被抛售，因此熊市中MA的边际定价效应更强。
第三，A 股熊市中常伴随 “救市” 政策与 “政策底” 预期，部分低共识股票可能因政策刺激出现短期反弹，
削弱了高低共识组合的收益差异，这是熊市分组检验显著性下降的重要原因。


